\chapter{थेल्स प्रमेय}\label{ch:g9:thales}

किसी \omterm{def:g8:solids:def}{पिरामिड} की ऊँचाई उस पर चढ़े बिना कैसे नापी जाए? मिलेतुस के थेल्स
ने उसकी छाया की तुलना एक छड़ी की छाया से की। उनकी प्रमेय --- \omterm{def:g6:lines:perp}{समांतर
रेखाएँ} \omterm{def:g6:lines:objects}{रेखाखंडों} को \omterm{def:g9:linfunc:linear}{समानुपाती} टुकड़ों में काटती हैं --- नमूनों,
नक़्शों और विवर्धनों की \omterm{def:g7:prop:scale}{मापनी} का गणितीय दिल है, और नाम वाली सबसे पुरानी
प्रमेयों में से एक।

\section{प्रमेय}

\begin{theorem}[थेल्स]\label{thm:g9:thales:direct}
मान लो दो रेखाएँ किसी बिंदु $A$ पर मिलती हैं। पहली रेखा पर $M$ लो और
दूसरी पर $N$; पहली पर $M$ से आगे $B$ है और दूसरी पर $N$ से आगे $C$।
अगर रेखाएँ $(MN)$ और $(BC)$ \emph{\omterm{def:g6:lines:perp}{समांतर}} हों, तो त्रिभुजों $AMN$ और
$ABC$ की भुजाएँ \omterm{def:g9:linfunc:linear}{समानुपाती} होती हैं:
\[
\frac{AM}{AB} = \frac{AN}{AC} = \frac{MN}{BC}.
\]
\end{theorem}
\admitted

\begin{omfigure}
\begin{tikzpicture}[scale=1.05]
  \coordinate (A) at (0,0);
  \coordinate (B) at (5.0,0.6);
  \coordinate (C) at (3.2,2.8);
  \coordinate (M) at ($(A)!0.55!(B)$);
  \coordinate (N) at ($(A)!0.55!(C)$);
  \draw[thick] (A) -- (B) -- (C) -- cycle;
  \draw[omThm, very thick] (M) -- (N);
  \draw[omThm, very thick] (B) -- (C);
  \fill (A) circle (1.5pt) node[left, font=\small] {$A$};
  \fill (B) circle (1.5pt) node[right, font=\small] {$B$};
  \fill (C) circle (1.5pt) node[above, font=\small] {$C$};
  \fill (M) circle (1.5pt) node[below, font=\small] {$M$};
  \fill (N) circle (1.5pt) node[above left, font=\small] {$N$};
\end{tikzpicture}\hspace{1.1cm}%
\begin{tikzpicture}[scale=1.05]
  \coordinate (A) at (0,0);
  \coordinate (B) at (2.6,-1.3);
  \coordinate (C) at (1.4,-2.2);
  \coordinate (M) at ($(A)!-0.62!(B)$);
  \coordinate (N) at ($(A)!-0.62!(C)$);
  \draw[thick] (M) -- (B);
  \draw[thick] (N) -- (C);
  \draw[omThm, very thick] (M) -- (N);
  \draw[omThm, very thick] (B) -- (C);
  \fill (A) circle (1.5pt) node[right=2pt, font=\small] {$A$};
  \fill (B) circle (1.5pt) node[right, font=\small] {$B$};
  \fill (C) circle (1.5pt) node[below, font=\small] {$C$};
  \fill (M) circle (1.5pt) node[left, font=\small] {$M$};
  \fill (N) circle (1.5pt) node[above, font=\small] {$N$};
\end{tikzpicture}

{\small थेल्स के दो विन्यास: एक के भीतर एक त्रिभुज (बाएँ) और
``तितली'', जिसमें $M$ तथा $N$ बिंदु $A$ के दूसरी ओर बैठते हैं (दाएँ)।
दोनों में $(MN) \parallel (BC)$ है और तीनों अनुपात बराबर हैं।}
\end{omfigure}

\begin{method}[थेल्स से कोई लंबाई निकालना]\label{met:g9:thales:compute}
\begin{enumerate}
  \item वह बिंदु $A$ पहचानो जहाँ दोनों रेखाएँ कटती हैं, और दोनों
        \omterm{def:g6:lines:perp}{समांतर रेखाएँ} भी; फिर दोनों त्रिभुजों के नाम लिखो;
  \item तीनों बराबर अनुपात लिखो, हमेशा ``छोटा त्रिभुज बटा बड़ा
        त्रिभुज'';
  \item उन्हीं दो अनुपातों की बराबरी रखो जिनमें तीनों ज्ञात लंबाइयाँ
        और अज्ञात लंबाई आती हैं;
  \item बने हुए \omterm{def:g7:prop:table}{समानुपात} को तिरछे नियम से हल करो।
\end{enumerate}
\end{method}

\begin{example}\label{ex:g9:thales:compute}
ऊपर बाएँ वाले विन्यास में मान लो $AM = 3$, $AB = 5$, $AN = 2.4$ और
$MN = 3.3$ हैं; $AC$ तथा $BC$ निकालते हैं। थेल्स देती है
\[
\frac{AM}{AB} = \frac{AN}{AC} = \frac{MN}{BC},
\qquad\text{यानी}\qquad
\frac35 = \frac{2.4}{AC} = \frac{3.3}{BC}.
\]
$\frac35 = \frac{2.4}{AC}$ से, तिरछा गुणा करने पर:
$3 \times AC = 5 \times 2.4 = 12$, यानी $AC = 4$। और
$\frac35 = \frac{3.3}{BC}$ से: $3 \times BC = 16.5$, यानी $BC = 5.5$।
\end{example}

\begin{example}[रोज़मर्रा में थेल्स]\label{ex:g9:thales:stick}
$1$ m ऊँची खड़ी छड़ी की छाया $1.5$ m की है, और उसी पल किसी पेड़ की
छाया $12$ m की। सूरज की किरणें \omterm{def:g6:lines:perp}{समांतर} हैं, इसलिए छड़ी-और-छाया तथा
पेड़-और-छाया वाले त्रिभुज थेल्स विन्यास में बैठते हैं:
\[
\frac{\text{पेड़ की ऊँचाई}}{1} = \frac{12}{1.5},
\qquad\text{यानी पेड़ } 8 \text{ m ऊँचा है।}
\]
\end{example}

\section{विलोम}

\begin{theorem}[थेल्स की विलोम]\label{thm:g9:thales:converse}
बिंदु वैसे ही रखे हों जैसे \cref{thm:g9:thales:direct} में ($M$ रेखा
$(AB)$ पर, $N$ रेखा $(AC)$ पर, दोनों रेखाओं पर एक ही क्रम में): अगर
\[
\frac{AM}{AB} = \frac{AN}{AC},
\]
हो, तो रेखाएँ $(MN)$ और $(BC)$ \omterm{def:g6:lines:perp}{समांतर} हैं।
\end{theorem}
\admitted

\begin{method}[समांतरता सिद्ध करना या ख़ारिज करना]\label{met:g9:thales:converse}
\begin{enumerate}
  \item दोनों अनुपात $\dfrac{AM}{AB}$ और $\dfrac{AN}{AC}$ अलग-अलग
        निकालो (भिन्नों के रूप में, \omterm{def:g6:decimals:rounding}{गोलन} के रूप में नहीं);
  \item अगर वे बराबर हों \emph{और} बिंदु दोनों रेखाओं पर एक ही क्रम
        में हों, तो रेखाएँ \omterm{def:g6:lines:perp}{समांतर} हैं (थेल्स की विलोम);
  \item अगर वे अलग हों, तो रेखाएँ \omterm{def:g6:lines:perp}{समांतर} \emph{नहीं} हैं --- क्योंकि
        अगर होतीं, तो थेल्स प्रमेय दोनों अनुपातों को बराबर कर देती।
\end{enumerate}
\end{method}

\begin{example}\label{ex:g9:thales:converse}
पहली रेखा पर $AM = 4$ और $AB = 10$; दूसरी पर $AN = 6$ और $AC = 15$।
तब
\[
\frac{AM}{AB} = \frac{4}{10} = \frac25,
\qquad
\frac{AN}{AC} = \frac{6}{15} = \frac25 :
\]
अनुपात बराबर, क्रम भी वही: $(MN) \parallel (BC)$।

इसके बजाय $AC = 14$ हो: $\frac{6}{14} = \frac37 \neq \frac25$, यानी
रेखाएँ \omterm{def:g6:lines:perp}{समांतर} नहीं हैं।
\end{example}

\section{विवर्धन और संकुचन}

\begin{definition}[आकृति को घटाना-बढ़ाना]\label{def:g9:thales:scaling}
किसी आकृति को \emph{मापक गुणक}\index{मापक गुणक} $k > 0$ से बढ़ाने या
घटाने का मतलब है उसकी \emph{सारी} लंबाइयों को $k$ से गुणा कर देना:
$k > 1$ पर \emph{विवर्धन}, और $k < 1$ पर \emph{संकुचन}। थेल्स के
विन्यास में त्रिभुज $AMN$, त्रिभुज $ABC$ का $k = \frac{AM}{AB}$ गुणक
वाला संकुचन है।
\end{definition}

\begin{proposition}[कोणों और क्षेत्रफलों पर असर]\label{prop:g9:thales:scaling}
गुणक $k$ से घटाने-बढ़ाने पर कोण और आकृतियों का आकार वही रहते हैं, हर
लंबाई $k$ से गुणा हो जाती है, और हर \emph{\omterm{def:g6:measure:area}{क्षेत्रफल}} $k^2$ से।
\end{proposition}
\admitted

\begin{omfigure}
\begin{tikzpicture}[scale=0.72]
  \draw[thick, omDef, fill=omDef!8] (0,0) rectangle (2,1.2);
  \draw[thick, omThm, fill=omThm!8] (3.4,0) rectangle (7.4,2.4);
  \node[font=\small, omDef] at (1,-0.45) {$2 \times 1.2$};
  \node[font=\small, omThm] at (5.4,-0.45) {$4 \times 2.4$};
  \node[font=\small] at (1,0.6) {क्षेत्रफल $2.4$};
  \node[font=\small] at (5.4,1.2) {क्षेत्रफल $9.6$};
\end{tikzpicture}

{\small लंबाइयाँ दुगनी करने ($k = 2$) पर \omterm{def:g6:measure:area}{क्षेत्रफल} $k^2 = 4$ गुना हो
जाता है: छोटे आयत की चार नक़लें बड़े को पूरा ढक देती हैं।}
\end{omfigure}

\begin{example}
$10 \times 15$ cm की कोई तस्वीर $k = 3$ \omterm{def:g9:thales:scaling}{मापक गुणक} से बड़ी की जाती है:
छपाई $30 \times 45$ cm की निकलती है। उसका \omterm{def:g6:measure:area}{क्षेत्रफल} $150$ cm$^2$ से
बढ़कर $150 \times 9 = 1350$ cm$^2$ हो जाता है --- $9$ गुना, $3$ गुना
नहीं।
\end{example}

\section{अभ्यास}

\begin{exercise}[$\star$]\label{exo:g9:thales:1}
रेखाएँ $(MN)$ और $(BC)$ \omterm{def:g6:lines:perp}{समांतर} हैं, $M$ बिंदु $[AB]$ पर है और $N$
बिंदु $[AC]$ पर; $AM = 2$, $AB = 6$, $AN = 3$, $BC = 9$। $AC$ और $MN$
निकालो।
\end{exercise}

\begin{exercise}[$\star$]\label{exo:g9:thales:2}
वही विन्यास: $AM = 5$, $MB = 3$ (ध्यान रहे: $MB$, $AB$ नहीं!),
$AN = 4$। $AC$ निकालो।
\end{exercise}

\begin{exercise}[$\star$]\label{exo:g9:thales:3}
तितली विन्यास में $A$ बिंदु $M$ और $B$ के बीच है, तथा $N$ और $C$ के
बीच भी, और $(MN) \parallel (BC)$ है; $AM = 3$, $AB = 7.5$, $AN = 2$,
$MN = 2.6$। $AC$ और $BC$ निकालो।
\end{exercise}

\begin{exercise}[$\star$]\label{exo:g9:thales:4}
$M$ बिंदु $[AB]$ पर है, $AM = 6$, $AB = 8$; $N$ बिंदु $[AC]$ पर है,
$AN = 9$, $AC = 12$। क्या रेखाएँ $(MN)$ और $(BC)$ \omterm{def:g6:lines:perp}{समांतर} हैं?
\end{exercise}

\begin{exercise}[$\star\star$]\label{exo:g9:thales:5}
ऊपर जैसा ही, पर $AM = 4$, $AB = 6$, $AN = 5$, $AC = 8$ के साथ। क्या
$(MN)$ और $(BC)$ \omterm{def:g6:lines:perp}{समांतर} हैं? ध्यान से कारण बताओ।
\end{exercise}

\begin{exercise}[$\star\star$]\label{exo:g9:thales:6}
$1.8$ m लंबा कोई व्यक्ति सड़क के खंभे की जड़ से $2$ m दूर खड़ा है;
उसकी छाया $3$ m की है। खंभे, व्यक्ति और छाया के सिरे से बनने वाला
थेल्स विन्यास बनाओ, और खंभे की ऊँचाई निकालो।
\end{exercise}

\begin{exercise}[$\star\star$]\label{exo:g9:thales:7}
किसी नक़्शे की \omterm{def:g7:prop:scale}{मापनी} $1 : 25\,000$ है (नक़्शे की लंबाइयाँ असली
लंबाइयों को $k = \frac{1}{25000}$ से गुणा करने पर मिलती हैं)।
\begin{enumerate}
  \item नक़्शे पर दो गाँव $6.8$ cm दूर हैं। असली दूरी कितनी है, km
        में?
  \item नक़्शे पर किसी जंगल का \omterm{def:g6:measure:area}{क्षेत्रफल} $8$ cm$^2$ है। उसका असली
        \omterm{def:g6:measure:area}{क्षेत्रफल} कितना है, km$^2$ में?
\end{enumerate}
\end{exercise}

\begin{exercise}[$\star\star$]\label{exo:g9:thales:8}
त्रिभुज $ABC$ में $AB = 12$, $AC = 15$, $BC = 18$ हैं। $[AB]$ पर बिंदु
$M$ ऐसा है कि $AM = 8$, और $M$ से होकर $(BC)$ के \omterm{def:g6:lines:perp}{समांतर} खींची गई रेखा
$[AC]$ को $N$ पर काटती है।
\begin{enumerate}
  \item $AN$ और $MN$ निकालो।
  \item $ABC$ से $AMN$ तक का \omterm{def:g9:thales:scaling}{मापक गुणक} क्या है? दोनों त्रिभुजों के
        परिमापों की तुलना करो।
\end{enumerate}
\end{exercise}

\begin{exercise}[$\star\star\star$]\label{exo:g9:thales:9}
मान लो $ABCD$ एक समलंब है जिसमें $(AB) \parallel (CD)$, $AB = 4$ और
$CD = 6$ हैं, और जिसके विकर्ण $[AC]$ तथा $[BD]$ बिंदु $O$ पर मिलते
हैं। $O$ के चारों ओर तितली विन्यास में थेल्स लगाकर अनुपात
$\dfrac{OA}{OC}$ निकालो, और दिखाओ कि $O$ दोनों विकर्णों को एक ही
अनुपात में काटता है।
\end{exercise}

\section{समस्या: बिना निशान की पट्टी, सूक्ष्मछिद्र कैमरा, और समलंब के
तीन माध्य}

\begin{problem}[{सप्ताहांत समस्या --- काम पर थेल्स: किसी भी रेखाखंड
को बराबर टुकड़ों में बाँटना, जूते के डिब्बे से सूरज नापना, और वह
समलंब जिसमें तीनों मशहूर माध्य आ मिलते हैं}]
\label{pb:g9:thales:1}
थेल्स प्रमेय \emph{किरणों} का गणित है: सूरज की किरणें, छोटे-से छेद से
गुज़रती रोशनी की किरणें, किसी \omterm{def:g5:solids:def}{शीर्ष} से निकलती पेंसिल की किरणें। यह
समस्या उसे तीन तरह से बरतती है --- रचना के औज़ार की तरह
(\emph{किसी भी} लंबाई का \omterm{def:g6:lines:objects}{रेखाखंड}, यहाँ तक कि $\sqrt2$ लंबाई का भी,
बिलकुल बराबर टुकड़ों में बाँटना), नापने के यंत्र की तरह (जूते के
डिब्बे और एक सिक्के से सूरज का \omterm{def:g4:geometry:circle}{व्यास}), और \cref{exo:g9:thales:9} वाले
समलंब पर आवर्धक शीशे की तरह, जहाँ \omterm{def:g6:lines:perp}{समांतर}, गुणोत्तर और हरात्मक ---
तीनों माध्य एक ही चित्र में आ जुटते हैं।

\medskip
\noindent\textbf{भाग I --- बिना निशान की पट्टी से बाँटना।}
\begin{enumerate}
  \item $7$ cm का \omterm{def:g6:lines:objects}{रेखाखंड} $[AB]$ खींचो। परकार और बिना निशान की
        पट्टी से उसे तीन बराबर टुकड़ों में काटने के लिए: $A$ से कोई
        किरण खींचो (जो $B$ से होकर न जाए); उस पर परकार से तीन बराबर
        लंबाइयाँ नापते जाओ, जिससे बिंदु $P_1$, $P_2$, $P_3$ मिलें;
        $P_3$ को $B$ से जोड़ो; फिर $P_1$ और $P_2$ से होकर
        $(P_3 B)$ के \omterm{def:g6:lines:perp}{समांतर रेखाएँ} खींचो। यह रचना करके निशान लगाओ कि
        वे \omterm{def:g6:lines:perp}{समांतर रेखाएँ} $[AB]$ को कहाँ काटती हैं।
  \item \cref{thm:g9:thales:direct} से कारण दो कि दोनों निशान लगे
        बिंदु $[AB]$ को ठीक उसके तिहाई बिंदुओं पर काटते हैं।
  \item इसी विधि को बदलकर किसी \omterm{def:g6:lines:objects}{रेखाखंड} को $5$ बराबर टुकड़ों में
        बाँटो, फिर किसी दिए हुए \omterm{def:g6:lines:objects}{रेखाखंड} का $\frac35$ बनाओ।
  \item $[AB]$ का वह बिंदु $M$ बनाओ जिसके लिए
        $\frac{AM}{MB} = \frac23$ हो (यानी $M$, $A$ से रास्ते के
        $\frac25$ पर)। किरण पर कितने बराबर क़दम चाहिए, और कौन-सा
        बिंदु $B$ से जोड़ा जाएगा?
  \item यह रचना नापने वाली पट्टी से नाप लेने से बेहतर क्यों है?
        $\sqrt2$ लंबाई का \omterm{def:g6:lines:objects}{रेखाखंड} सोचो (एक इकाई वर्ग का विकर्ण, जो
        \cref{pb:g9:sqrt:1} के अनुसार अपरिमेय है): नापने से क्या
        मिलता, और थेल्स क्या देती है?
\end{enumerate}

\medskip
\noindent\textbf{भाग II --- सूक्ष्मछिद्र कैमरा।}
किसी बंद डिब्बे के एक \omterm{def:g5:solids:def}{फलक} में सुई से छेद करो: सामने वाले \omterm{def:g5:solids:def}{फलक} पर
दुनिया की उलटी तस्वीर उभर आती है। वस्तु का कोई बिंदु, छेद, और तस्वीर
का बिंदु --- तीनों एक ही रेखा पर होते हैं, क्योंकि रोशनी सीधी चलती
है --- इसलिए वस्तु (ऊँचाई $H$, छेद के सामने $D$ दूरी पर) और उसकी
तस्वीर (ऊँचाई $h$, छेद के पीछे $d$ गहराई पर पिछली दीवार पर) थेल्स के
तितली विन्यास में बैठती हैं।
\begin{enumerate}[resume]
  \item छेद के चारों ओर तितली में थेल्स लगाकर सूक्ष्मछिद्र सूत्र
        $h = H \times \frac dD$ निकालो।
  \item $12$ m ऊँचा कोई पेड़ $20$ cm गहरे जूते के डिब्बे से $30$ m
        दूर खड़ा है। उसकी तस्वीर कितनी ऊँची बनेगी, और सीधी या उलटी?
  \item डिब्बे को सूरज की ओर करो: छेद के $d = 1$ m पीछे सूरज की
        तस्वीर क़रीब $9$ mm \omterm{def:g4:geometry:circle}{व्यास} की एक चकती बनती है। सूरज की दूरी
        $D = 1.5 \times 10^{11}$ m मानकर उसका \omterm{def:g4:geometry:circle}{व्यास} \omterm{def:g9:fractions:scientific}{वैज्ञानिक लेखन}
        में निकालो (\cref{pb:g9:fractions:1} वाला लेखन काम पर)। असली
        मान $1.39 \times 10^9$ m से तुलना करो।
  \item अब चाँद: \omterm{def:g4:geometry:circle}{व्यास} $3.5 \times 10^6$ m, दूरी
        $3.8 \times 10^8$ m। चाँद के लिए और सूरज के लिए
        $\frac{\text{व्यास}}{\text{दूरी}}$ अनुपात निकालो। दोनों
        संख्याएँ किस भाग्यशाली संयोग की ओर इशारा करती हैं --- और वह
        कौन-सा शानदार नज़ारा मुमकिन बनाता है?
  \item एक-दो वाक्यों में: सूक्ष्मछिद्र की तस्वीर उलटी क्यों होती
        है, और छेद बड़ा करने पर वह चमकीली तो हो जाती है पर धुँधली
        भी क्यों पड़ जाती है?
\end{enumerate}

\medskip
\noindent\textbf{भाग III --- तीन माध्यों वाला समलंब।}
$ABCD$ वही समलंब है जो \cref{exo:g9:thales:9} में था:
$(AB) \parallel (CD)$, $AB = 4$, $CD = 6$, और विकर्ण $O$ पर कटते हैं,
जहाँ $\frac{OA}{OC} = \frac{OB}{OD} = \frac{4}{6} = \frac23$ है।
\begin{enumerate}[resume]
  \item तितली का अभ्यास: दो रेखाएँ $O$ पर कटती हैं और
        $(MN) \parallel (PQ)$ तितली की स्थिति में हैं; $OM = 3$,
        $OP = 7.5$, $MN = 4$ और $OQ = 6$। $PQ$ तथा $ON$ निकालो।
  \item समझाओ कि $\frac12$ अनुपात के साथ लगाई गई थेल्स में
        \cref{thm:g8:midpoints:direct} की मध्यबिंदु प्रमेय एक ख़ास
        हालत की तरह समाई हुई है।
  \item $O$ से होकर आधारों के \omterm{def:g6:lines:perp}{समांतर} रेखा खींचो; वह $[AD]$ को $E$ पर
        मिलती है। त्रिभुज $ACD$ में काम करते हुए
        ($\frac{AO}{AC} = \frac25$ पर ध्यान दो) $EO$ निकालो।
  \item त्रिभुज $BDC$ में वैसी ही दर्पण-गणना से टुकड़ा $OF$ ($F$
        बिंदु $[BC]$ पर) उतना ही लंबा निकलता है। इससे नतीजा निकालो
        कि
        \[
        EF = 2 \times \frac{4 \times 6}{4 + 6} = 4.8 :
        \]
        यानी विकर्णों के कटान-बिंदु से गुज़रती \omterm{def:g6:lines:perp}{समांतर} रेखा आधारों का
        \emph{\omterm{pb:g8:speed:1}{हरात्मक माध्य}} नापती है --- \cref{pb:g8:speed:1} वाला
        आने-जाने का माध्य, जो यहाँ शुद्ध ज्यामिति में फिर से आ खड़ा
        होता है।
  \item भव्य समापन: आधारों $4$ और $6$ के तीनों शास्त्रीय माध्य
        निकालो --- \omterm{def:g6:lines:perp}{समांतर} $\frac{4+6}{2}$, गुणोत्तर
        $\sqrt{4 \times 6}$, हरात्मक
        $\frac{2 \times 4 \times 6}{4 + 6}$ --- और उन्हें क्रम में
        लगाओ। हर एक समलंब में बसता है: \omterm{def:g6:lines:perp}{समांतर} माध्य वह मध्य रेखा है
        जो असमांतर भुजाओं के मध्यबिंदुओं को जोड़ती है
        (\cref{pb:g7:areas:1}), \omterm{pb:g8:speed:1}{हरात्मक माध्य} तुम्हारा \omterm{def:g6:lines:objects}{रेखाखंड} $EF$
        है, और \omterm{pb:g8:cosine:1}{गुणोत्तर माध्य} वह \omterm{def:g6:lines:perp}{समांतर} रेखा है जो समलंब को दो
        \emph{समरूप} समलंबों में काट देती है (उसका मान कैलकुलेटर से
        जाँचो; प्रमाण एक प्यारी अतिरिक्त चुनौती है)। इन तीनों के बीच
        की असमिका इस किताब में कहाँ सिद्ध हुई है
        (\cref{pb:g8:cosine:1}, \cref{pb:g8:speed:1})?

\end{enumerate}
\end{problem}
